\documentclass[11pt,a4paper]{article}
\usepackage[margin=1in]{geometry}
\usepackage[T1]{fontenc}
\usepackage[utf8]{inputenc}
\usepackage{booktabs}
\usepackage{longtable}
\begin{document}
\section*{Tier 1 Visual Adjudication}
This report converts the Tier 1 visual evidence into provisional research decisions. It does not merge signs; it separates visual allograph candidates from distributional functional clusters.
\begin{longtable}{p{0.10\textwidth}p{0.20\textwidth}p{0.28\textwidth}p{0.28\textwidth}}
\toprule
\textbf{Pair} & \textbf{Shape assessment} & \textbf{Provisional verdict} & \textbf{Next action} \\
\midrule
\endhead
817/861 & Moderate family resemblance: both are enclosed signs with angular internal marking, but outline and internal geometry differ. & Keep distinct; functional-cluster candidate, not a merge candidate. & Check artifact photos for carving distortion and independent occurrences. \\
817/820 & Moderate family resemblance: both are enclosed signs, but 817 has a simple chevron-like interior and 820 has a divided/crossed interior. & Keep distinct; functional-cluster candidate, not a merge candidate. & Check artifact photos for systematic graphic alternation or catalog separation. \\
705/706 & Strong visual resemblance: same open-U/long-stem family, with 706 adding or emphasizing an inner nested stroke. & Variant-link candidate; no merge until artifact-level review. & Prioritize CISI image checks for crowding, carving depth, and object-level independence. \\
692/920 & Weak visual resemblance: 692 is crossed/X-shaped, while 920 is a curved single-stroke form. & Reject visual-allograph hypothesis at catalog level; keep as functional-cluster candidate. & Use as a control case for distributional similarity without visual similarity. \\
\bottomrule
\end{longtable}
\end{document}
